% Blatt 7, A4. Aufgabenstellung und Skizze: aufgaben/Blatt07.pdf, S. 2.

\subsection*{Geometrie und Modell}
Das L-förmige Gefäss ist bis zur Oberkante mit Wasser gefüllt. Die Skizze
zeigt die Höhe \(3d\) und den Abstand \(2d\) von der Oberkante bis zur Stufe.
Die quadratische Fläche A ist die \emph{waagerechte Deckfläche} des niedrigen
Arms, mit Wasser darunter. Sie liegt somit durchgehend in der Tiefe \(2d\).
Die ebenfalls quadratische Fläche B ist die \emph{senkrechte rechte
Stirnfläche} des Arms und reicht von der Tiefe \(2d\) bis \(3d\). Beide
haben den Flächeninhalt \(d^2\).

Wir betrachten ruhendes Wasser konstanter Dichte \(\rho\) in einem homogenen
Schwerefeld \(g\). Der Druck an der offenen Wasseroberfläche sei \(p_0\).
Zuerst bestimmen wir die Druckkräfte, die das Wasser auf die jeweiligen
Flächen ausübt. Da \(d\), \(\rho\), \(g\) und \(p_0\) nicht numerisch
vorgegeben sind, bleiben die Ergebnisse allgemein.

\subsection*{a) Kraft auf die waagerechte Fläche A}
Auf A ist die Tiefe und damit der Druck überall gleich. Das Wasser drückt von
unten nach oben. Die reine Wasserkraft ist
\begin{equation}
  F_{A,\mathrm{Wasser}}=(p_0+2\rho gd)d^2
  \quad\text{nach oben}.
\end{equation}
Ist die Oberseite von A Luft mit dem Druck \(p_0\) ausgesetzt, drückt diese
entgegengesetzt. Die resultierende \emph{Druckkraft} auf das Wandstück ist
\begin{ergebnis}[a) Kraft aus der Druckdifferenz]
\begin{equation*}
  F_{A,\mathrm{netto}}=2\rho gd^3\quad\text{nach oben}.
\end{equation*}
\end{ergebnis}

\subsection*{b) Kraft auf die senkrechte Fläche B}
Auf B nimmt der Druck mit der Tiefe zu. Ein horizontaler Streifen in der Tiefe
\(h\) hat die Höhe \(\mathrm dh\) und die Breite \(d\), also
\(\mathrm dA=d\,\mathrm dh\). Die reine Wasserkraft nach rechts ist daher
\begin{align}
  F_{B,\mathrm{Wasser}}
  &=\int_{2d}^{3d}(p_0+\rho gh)\,d\,\mathrm dh\\
  &=p_0d^2+\rho gd\left[\frac{h^2}{2}\right]_{2d}^{3d}
    =p_0d^2+\frac52\rho gd^3.
\end{align}
Wenn aussen gleichmässig der Druck \(p_0\) wirkt, ist die resultierende
Druckkraft auf das Wandstück
\begin{ergebnis}[b) Kraft aus der Druckdifferenz]
\begin{equation*}
  F_{B,\mathrm{netto}}=\frac52\rho gd^3\quad\text{nach rechts}.
\end{equation*}
\end{ergebnis}
Ohne gleichen Aussendruck gelten diese Nettoergebnisse nicht unverändert:
Bei konstantem Aussendruck \(p_{\mathrm{a}}\) ersetzt man in beiden
Wasserkräften den Term \(p_0d^2\) durch \((p_0-p_{\mathrm{a}})d^2\).
Die Kräfte von Befestigungen und das Eigengewicht der Wandstücke sind dabei
nicht enthalten.

\begin{kontrolle}[Plausibilität]
Für B ist die mittlere Tiefe \(\tfrac52d\). Wegen des linear mit der Tiefe
wachsenden Schweredrucks ist der mittlere Überdruck dort
\(\tfrac52\rho gd\); mal \(d^2\) bestätigt er das Integral. Die
Kraftformeln haben die Einheit \(\si{\newton}\). Im Grenzfall \(g=0\)
verschwindet bei gleichem Umgebungsdruck die resultierende Druckkraft.
\end{kontrolle}

\begin{merksatz}
Bei räumlich veränderlichem Druck addiert man die lokalen Druckkräfte. Für
eine ebene Fläche ist dabei auch die Druckbelastung der Rückseite zu prüfen.
\end{merksatz}
